% EABC Final Paper: Character-Theoretic Diagonalization of mod-420 Channel Asymmetry
% Compile: pdflatex eabc_paper_final.tex
\documentclass[11pt,a4paper]{article}

\usepackage[utf8]{inputenc}
\usepackage[T1]{fontenc}
\usepackage{amsmath,amssymb,amsthm,mathtools}
\usepackage{booktabs}
\usepackage[a4paper,margin=2.5cm]{geometry}
\usepackage{microtype}
\usepackage{hyperref}

\title{Character-Theoretic Diagonalization of the\\
Modulo-$420$ Channel Asymmetry in Prime Quadruplets}
\author{Author Name(s) \and Affiliation(s)}
\date{\today}

\newtheorem{theorem}{Theorem}
\newtheorem{lemma}[theorem]{Lemma}
\newtheorem{corollary}[theorem]{Corollary}
\newtheorem{proposition}[theorem]{Proposition}
\newtheorem{hypothesis}{Hypothesis}
\newtheorem{remark}{Remark}

\newcommand{\R}{\mathcal{R}}
\newcommand{\E}{\mathbb{E}}
\newcommand{\Var}{\mathrm{Var}}

\begin{document}

\maketitle

\begin{abstract}
We study the occupancy asymmetry on the six observed prime-quadruplet
half-channels $\R=\{11,101,191,221,311,401\}$ modulo $420$.
Three methodological levels structure the analysis.
\textbf{Level~1 (exact algebra):}
the primitive cubic Dirichlet character $\chi_3\bmod 7$ coincides
pointwise with the $k{=}2$ discrete Fourier mode on $\R$
($\chi_3|_{\R}=\psi_2$), hence $F_2=\langle\delta,\chi_3\rangle$ for the
observed bias vector $\delta$; the bias admits an exact orthogonal
decomposition into Dirichlet character components modulo $420$.
\textbf{Level~2 (numerical, conditional on $H$):}
under a multinomial equidistribution hypothesis $H$, the $\chi_3$-projection
scales as $|\langle\delta_N,\chi_3\rangle|\sim 1/(6\sqrt{N})$ with
$z$-scores $\approx 1$ at $N_{\max}\in\{10^8,3\cdot 10^8,10^9\}$---pure
fluctuation, no persistent Chebyshev-type drift.
\textbf{Level~3 (HL corroboration):}
a Hardy--Littlewood singular-series analysis yields \emph{identical}
coefficients on all six channels in $\R$, predicting asymptotic
equidistribution $1/6$; the observed asymmetry is statistically
consistent with finite-sample fluctuations ($\chi^2\approx 4.34$,
$p\approx 0.50$).
Analytic proof of conditional equidistribution remains open.
\end{abstract}

%=============================================================================
\section{Three-Level Framework}
\label{sec:three-levels}
%=============================================================================

The present work separates three epistemically distinct layers.
Only Level~1 consists of unconditional exact statements; Level~2 is
probabilistic and conditional on Hypothesis~$H$; Level~3 remains open.

\subsection{Level~1: Exact Algebraic Statements}

These results hold for the \emph{fixed finite channel set} $\R$ and the
\emph{observed sample} bias $\delta$ (or $\delta_N$ at any finite count $N$).
They do not assert asymptotic bias in $N$.

\begin{enumerate}
  \item \textbf{Pointwise identity.}
        $\chi_3|_{\R}=\psi_2$ as vectors in $\mathbb{C}^6$, where
        $\psi_k(n):=e^{2\pi i kn/6}$ in the fixed order
        $(11,101,191,221,311,401)$.
  \item \textbf{DFT--character coincidence.}
        For the observed (centered) bias vector $\delta$,
        $F_2=\langle\delta,\chi_3\rangle$, with $F_k$ the $k$-th DFT
        coefficient on the six-point set.
  \item \textbf{Orthogonal decomposition.}
        $\delta$ lies in the Dirichlet character subspace mod $420$:
        numerically,
        $\delta = P_{\chi_3}\delta + P_{\overline{\chi_3}}\delta
        + P_{\chi_4}\delta + P_{\chi_4\chi_3}\delta + \varepsilon$
        with $\|\varepsilon\|^2 < 10^{-36}$ for the sample at $N=4767$.
\end{enumerate}

\subsection{Level~2: Numerically Supported Under Hypothesis $H$}

Under the multinomial null model $H$ (Section~\ref{sec:hypothesis-H}),
the $\chi_3$-coefficient is a fluctuation of typical size
$f(N)=1/(6\sqrt{N})$:
\[
  |\langle\delta_N,\chi_3\rangle| \sim \frac{1}{6\sqrt{N}},
  \qquad z \approx 1,
\]
at $N_{\max}\in\{10^8,3\cdot 10^8,10^9\}$.
The pair bias $B_{101,11}/\sqrt{N}$ shows no stable positive drift.
The data are compatible with multinomial fluctuations around the uniform
distribution, not with a growing structural bias.

\subsection{Level~3: HL Corroboration and Open Equidistribution}

Section~\ref{sec:hl-analysis} computes the Hardy--Littlewood singular
series for the gap constellation $(0,2,6,8)$ and shows that all six
channels in $\R$ carry identical local factors mod $2,3,4,5,7$.
The HL heuristic therefore predicts asymptotic equidistribution
$P(p\bmod 420=r\mid Q(p)\text{ prime})\to 1/6$.
The observed occupancy vector at $N=4767$ is statistically consistent
with this prediction ($\chi^2\approx 4.34$, $p\approx 0.50$).
What remains \emph{unproved} is an analytic equidistribution theorem
for $p\bmod 420$ conditional on $Q(p)$ being prime, including control
of arithmetic correlations between successive events.

%=============================================================================
\section{Setup and Notation}
%=============================================================================

Let $Q(p)=(p,p+2,p+6,p+8)$ denote a prime quadruplet with gap signature
$(2,4,2)$.
The six observed half-channels modulo $420$ are
\[
  \R = \{11,101,191,221,311,401\}.
\]
For $N$ counted quadruplet events with channel label $r_k\in\R$, write
$n_r(N)$ for occupancy counts and $K=\sum_r n_r(N)=N$.

The \emph{centered channel bias} (probability form) is
\[
  \delta_N(r) := \frac{n_r(N)}{N} - \frac{1}{6},
  \qquad r\in\R.
\]
For the reference sample $N=4767$ we write $\delta:=\delta_N$.
The inner product on $\mathbb{C}^{\R}$ is
\[
  \langle f,g\rangle := \frac{1}{6}\sum_{r\in\R} f_r\,\overline{g_r}.
\]
Dirichlet characters mod $420=4\cdot 3\cdot 5\cdot 7$ are written as
product characters $\chi=\chi_4\cdot\chi_5\cdot\chi_7^{(q)}\cdot\chi_7^{(3)}$,
with $\chi_4$ the primitive quadratic character mod $4$,
$\chi_5=(\cdot/5)$, $\chi_7^{(q)}=(\cdot/7)$, and $\chi_7^{(3)}=\chi_3$
the primitive \emph{cubic} character mod $7$
($\chi_3(\zeta)=\zeta$ for $\zeta\in\mathbb{F}_7^\times$,
$\omega=e^{2\pi i/3}$, $\chi_3(4)=1$, $\chi_3(3)=\omega$, $\chi_3(2)=\omega^2$).

%=============================================================================
\section{Character-Theoretic Diagonalization}
\label{sec:diagonalization}
%=============================================================================

We state only what is proved or numerically verified for the \emph{observed
finite sample}; coefficients do not grow with $N$ in the algebraic layer.

\begin{lemma}[Mod-$7$ cycle on $\R$]
\label{lem:mod7-cycle}
In the fixed channel order $(11,101,191,221,311,401)$,
\[
  r_n \bmod 7 = (4,\,3,\,2,\,4,\,3,\,2).
\]
\end{lemma}

\begin{proof}
Direct computation:
$11\equiv 4$, $101\equiv 3$, $191\equiv 2$,
$221\equiv 4$, $311\equiv 3$, $401\equiv 2 \pmod{7}$.
\end{proof}

\begin{lemma}[Pointwise identity $\chi_3=\psi_2$ on $\R$]
\label{lem:chi3-equals-psi2}
With $\omega=e^{2\pi i/3}$,
\[
  \chi_3(r_n) = \psi_2(n)
  \qquad\text{for all } n=0,\ldots,5.
\]
Hence $\chi_3|_{\R}$ and $\psi_2$ are \emph{identical} vectors in
$\mathbb{C}^6$ (not merely collinear).
\end{lemma}

\begin{proof}
By Lemma~\ref{lem:mod7-cycle},
$(\chi_3(r_n))_{n=0}^{5}=(1,\omega,\omega^2,1,\omega,\omega^2)$.
Since $\psi_2(n)=e^{2\pi i\cdot 2n/6}=e^{2\pi i n/3}$, the sequences agree.
\end{proof}

\begin{theorem}[Character-Theoretic Diagonalization]
\label{thm:character-diagonalization}
\emph{(For the observed finite sample on $\R$.)}
Let $(\delta_n)_{n=0}^{5}$ be $\delta$ in the fixed channel order and
\[
  F_k := \frac{1}{6}\sum_{n=0}^{5} \delta_n\, e^{-2\pi i kn/6}
  = \langle\delta,\psi_k\rangle
\]
the $k$-th DFT coefficient. Then:
\begin{enumerate}
  \item \textbf{Exact coincidence.}
        $F_2 = \langle\delta,\chi_3\rangle$.
        Since $\delta\in\mathbb{R}^6$, also $F_4=\overline{F_2}$.
  \item \textbf{Energy share.}
        The projection onto $\mathrm{span}\{\chi_3,\overline{\chi_3}\}$
        accounts for
        $|F_2|^2+|F_4|^2
        = |\langle\delta,\chi_3\rangle|^2
        + |\langle\delta,\overline{\chi_3}\rangle|^2
        = 0.5971\cdot\Var(\delta)$
        in the sample at $N=4767$.
  \item \textbf{Complete character decomposition.}
        With orthonormal character directions $\{\hat{e}_\chi\}$ on $\R$
        (Gram--Schmidt on character values),
        \[
          \delta = \sum_\chi c_\chi\,\hat{e}_\chi,
          \qquad c_\chi=\langle\delta,\hat{e}_\chi\rangle,
          \qquad \Var(\delta)=\sum_\chi|c_\chi|^2.
        \]
        Numerically,
        \[
          \delta = P_{\chi_3}\delta + P_{\overline{\chi_3}}\delta
                   + P_{\chi_4}\delta + P_{\chi_4\chi_3}\delta + \varepsilon,
          \qquad \|\varepsilon\|^2 < 10^{-36}.
        \]
        The residual after removing $\chi_3+\overline{\chi_3}$ is carried
        primarily by $\chi_4$ ($10.67\%$ of total variance) and
        $\chi_4\chi_3$ ($28.75\%$); $\chi_5$ vanishes since $r\equiv 1\pmod 5$
        on all channels.
\end{enumerate}
\end{theorem}

\begin{proof}
\emph{Part (1).}
Lemma~\ref{lem:chi3-equals-psi2} gives $\psi_2=\chi_3|_{\R}$, hence
$F_2=\langle\delta,\psi_2\rangle=\langle\delta,\chi_3\rangle$.
Reality of $\delta$ and $\psi_4=\overline{\psi_2}$ yield $F_4=\overline{F_2}$.

\emph{Part (2).}
Parseval on the six-point set; modes $k=2,4$ span the same two-dimensional
cubic subspace as $\{\chi_3,\overline{\chi_3}\}$.

\emph{Part (3).}
All $r\in\R$ are odd and $\not\equiv 0\pmod{5,7}$, so all characters are
well-defined.
On $\R$, $r\equiv 1\pmod 5$ and $r\equiv 2\pmod 3$, hence $\chi_5$ and
the quadratic character mod $3$ are constant and contribute no variance.
Gram--Schmidt reconstruction gives
$\Var(\varepsilon)=1.31\times 10^{-36}$.
On the $\chi_3=1$ subclass $\{11,221\}$, only $\chi_4$ separates the two
channels ($\chi_4(11)=-1$, $\chi_4(221)=+1$).
\end{proof}

\begin{remark}[Scope of the theorem]
Theorem~\ref{thm:character-diagonalization} is a \emph{diagonalization of
fluctuations in character coordinates} for the observed finite sample.
It identifies $\chi_3$ as the dominant Fourier direction on $\R$ by
\emph{algebraic} coincidence, not by asymptotic equidistribution.
It does \emph{not} claim that $|\langle\delta_N,\chi_3\rangle|$ grows with $N$,
nor that a Chebyshev-type bias persists.
\end{remark}

%=============================================================================
\section{Numerical Results}
\label{sec:numerical}
%=============================================================================

\subsection{Character decomposition at $N_{\max}=10^8$ ($K=4767$)}

\begin{table}[ht]
\centering
\caption{Orthonormal character projections of $\delta$ on $\R$;
variance shares in percent ($\Var(\delta)=2.53\times 10^{-5}$).}
\label{tab:projections}
\begin{tabular}{lrrr}
\toprule
Character $\chi$ & $|c_\chi|$ & $\arg(c_\chi)$ & Variance (\%) \\
\midrule
$1$ (trivial) & $0.00000000$ & --- & $0.0000$ \\
$\chi_3$ & $0.00274918$ & $-106.63^\circ$ & $29.8566$ \\
$\overline{\chi_3}$ & $0.00274918$ & $+106.63^\circ$ & $29.8566$ \\
$\chi_4$ & $0.00164324$ & $0.00^\circ$ & $10.6669$ \\
$\chi_4\chi_7^{(q)}$ & $0.00046972$ & $180.00^\circ$ & $0.8716$ \\
$\chi_4\chi_3$ & $0.00269768$ & $-30.00^\circ$ & $28.7484$ \\
\midrule
\textbf{Sum} & --- & --- & \textbf{100.0000} \\
\bottomrule
\end{tabular}
\end{table}

\begin{table}[ht]
\centering
\caption{Primitive character values on $\R$.}
\label{tab:char-values}
\begin{tabular}{rrrrrrr}
\toprule
$r$ & $r\bmod 4$ & $r\bmod 5$ & $r\bmod 7$ & $\chi_4$ & $\chi_5$ & $\chi_3$ \\
\midrule
$11$  & $3$ & $1$ & $4$ & $-1$ & $+1$ & $+1$ \\
$101$ & $1$ & $1$ & $3$ & $+1$ & $+1$ & $\omega$ \\
$191$ & $3$ & $1$ & $2$ & $-1$ & $+1$ & $\omega^2$ \\
$221$ & $1$ & $1$ & $4$ & $+1$ & $+1$ & $+1$ \\
$311$ & $3$ & $1$ & $3$ & $-1$ & $+1$ & $\omega$ \\
$401$ & $1$ & $1$ & $2$ & $+1$ & $+1$ & $\omega^2$ \\
\bottomrule
\end{tabular}
\end{table}

\subsection{Scaling across $10^8$, $3\cdot 10^8$, $10^9$}

\begin{table}[ht]
\centering
\caption{$\chi_3$-projection vs.\ theory $f(K)=1/(6\sqrt{K})$ under $H$
(data from \texttt{eabc\_analytic\_proof\_check.py};
multinomial null: $10\,000$ trials).}
\label{tab:chi3-scaling}
\begin{tabular}{rrrrrrr}
\toprule
$N_{\max}$ & $K$ & $|\langle\delta,\chi_3\rangle|$ & $f(K)$ & $z$ &
$|\langle\delta,\chi_3\rangle|\sqrt{K}$ & $p(\ge)$ \\
\midrule
$10^8$       & $4767$  & $0.002749$ & $0.002414$ & $1.14$ & $0.190$ & $26.8\%$ \\
$3\cdot10^8$ & $10972$ & $0.001568$ & $0.001591$ & $0.99$ & $0.164$ & $37.7\%$ \\
$10^9$       & $28387$ & $0.001043$ & $0.000989$ & $1.05$ & $0.176$ & $33.1\%$ \\
\bottomrule
\end{tabular}
\end{table}

\begin{table}[ht]
\centering
\caption{Pair bias $B_{101,11}=N_{101}-N_{11}$
(data from \texttt{eabc\_scaling\_test.py}).}
\label{tab:B-scaling}
\begin{tabular}{rrrr}
\toprule
$N_{\max}$ & $K$ & $B_{101,11}$ & $B/\sqrt{K}$ \\
\midrule
$10^8$       & $4767$  & $+66$ & $+0.956$ \\
$3\cdot10^8$ & $10972$ & $+51$ & $+0.487$ \\
$10^9$       & $28387$ & $+77$ & $+0.457$ \\
\bottomrule
\end{tabular}
\end{table}

The product $|\langle\delta,\chi_3\rangle|\sqrt{K}\approx 0.18\pm 0.01$ is
stable across three decades; all $z$-scores lie near $1$ (Rayleigh median
$\approx 1.18$ under $H$).
The ratio $B/\sqrt{K}$ decreases monotonically ($0.956\to 0.487\to 0.457$)
without stabilizing at a positive constant; the terminal null $p$-value is
$\approx 43\%$.

%=============================================================================
\section{Proof Attempt: Hypothesis $H$ and $f(N)=1/(6\sqrt{N})$}
\label{sec:hypothesis-H}
%=============================================================================

\begin{hypothesis}[$H$---Multinomial equidistribution on $\R$]
\label{hyp:H}
For $N$ quadruplet events, channel occupancy follows
\[
  (n_{11},n_{101},\ldots,n_{401}) \sim
  \mathrm{Mult}\!\left(N;\,\tfrac{1}{6},\ldots,\tfrac{1}{6}\right).
\]
Equivalently, channel indices $(I_k)_{k=1}^{N}$ are i.i.d.\ uniform on
$\{1,\ldots,6\}$.
\end{hypothesis}

\begin{remark}
$H$ is the asymptotic consequence of Hardy--Littlewood equidistribution
\emph{if} all six channels in $\R$ carry equal singular-series factors.
$H$ does not postulate Chebyshev drift, rank-one dynamics, or causal EABC
interpretation.
\end{remark}

\begin{proposition}[CLT scaling under $H$]
\label{prop:clt}
Under Hypothesis~\ref{hyp:H}, as $N\to\infty$:
\begin{enumerate}
  \item $\E[\langle\delta_N,\chi_3\rangle]=0$;
  \item $\Var(\langle\delta_N,\chi_3\rangle)=1/(36N)$;
  \item $\sqrt{N}\,\langle\delta_N,\chi_3\rangle
        \xRightarrow{d}\mathcal{N}_{\mathbb{C}}(0,1/36)$;
  \item $|\langle\delta_N,\chi_3\rangle|=O_p(1/\sqrt{N})$.
\end{enumerate}
\end{proposition}

\begin{proof}[Proof attempt (conditional on $H$)]
Since $\sum_{r\in\R}\chi_3(r)=0$, the mean cancels:
\[
  \langle\delta_N,\chi_3\rangle
  = \frac{1}{6N}\sum_{r\in\R} n_r\,\overline{\chi_3(r)}
  = \frac{S_N}{6N},
  \qquad
  S_N := \sum_{k=1}^{N}\overline{\chi_3(r_{I_k})}.
\]
Under $H$, $Y_k:=\overline{\chi_3(r_{I_k})}$ are i.i.d.\ with $\E[Y_k]=0$ and
$\E[|Y_k|^2]=1$, hence $\Var(S_N)=N$ and
$\Var(S_N/(6N))=1/(36N)$.
The multidimensional CLT gives
$\sqrt{N}\,\langle\delta_N,\chi_3\rangle\xRightarrow{d}\mathcal{N}_{\mathbb{C}}(0,1/36)$.
\end{proof}

\begin{corollary}[Explicit scaling function]
\label{cor:fN}
Under $H$, the natural scale is
\[
  \boxed{f(N)=\frac{1}{6\sqrt{N}},}
  \qquad
  \langle\delta_N,\chi_3\rangle
  \sim \mathcal{N}_{\mathbb{C}}\!\left(0,\frac{1}{36N}\right).
\]
The $z$-score $|\langle\delta_N,\chi_3\rangle|/f(N)$ is asymptotically
Rayleigh$(1)$-distributed (median $\approx 1.177$).
Expected magnitude:
$\E[|\langle\delta_N,\chi_3\rangle|]\approx\sqrt{\pi/(72N)}\approx 0.886\,f(N)$.
\end{corollary}

\paragraph{Proof gap.}
The CLT argument is complete \emph{conditional on $H$}.
Section~\ref{sec:hl-analysis} supplies arithmetic evidence that $H$
is the natural HL consequence on $\R$; an analytic proof of conditional
equidistribution (Hypothesis~H6) remains open.

%=============================================================================
\section{Hardy--Littlewood Equidistribution Analysis}
\label{sec:hl-analysis}
%=============================================================================

Let $Q=(0,2,6,8)$ denote the gap constellation for quadruplets
$(p,p+2,p+6,p+8)$ and
\[
  \pi_Q(x) := \#\{p\le x : p,p+2,p+6,p+8\text{ prime}\}.
\]
The Hardy--Littlewood conjecture \cite{hardylittlewood1923} (Bateman--Horn
form) predicts
\[
  \pi_Q(x) \;\sim\; C_Q \int_2^x \frac{dt}{\log^4 t},
  \qquad
  C_Q = 2\cdot \mathfrak{S}(Q),
\]
where $\mathfrak{S}(Q)=\prod_p w_p(Q)$ is the \emph{singular series}.
For prime $p$, let $\nu_p$ be the number of forbidden start classes
mod $p$ (classes $n$ with $n+h\equiv 0$ for some $h\in Q$). Then
\begin{equation}
  w_p(Q) = 1 - \frac{\nu_p}{(p-1)^4}.
  \label{eq:wp}
\end{equation}
A partial product to $p\le 2\cdot 10^5$ gives
$\mathfrak{S}(Q)\approx 0.858$.

\paragraph{Conditional distribution mod $q$.}
For modulus $q$, let $\omega_q$ be the number of admissible start classes
($n\bmod q$ with all four positions $\not\equiv 0$).
Under HL equidistribution on admissible classes, $p\equiv r\bmod q$
carries relative weight $1/\omega_q$ if admissible, else $0$.
The total Hardy--Littlewood weight of channel $r\in\R$ is
\begin{equation}
  W_{\mathrm{HL}}(r) \;=\;
  \prod_{q\in\{2,3,4,5,7\}} \frac{1}{\omega_q}
  \cdot \prod_{p\in\{3,5,7\}} w_p(Q),
  \label{eq:channel-weight}
\end{equation}
provided $r$ is admissible at every factor.

\begin{lemma}[$\R$ is complete mod $420$]
\label{lem:R-complete}
$\omega_{420}=6$, and the only admissible residue classes are exactly
$\R=\{11,101,191,221,311,401\}$.
\end{lemma}

\begin{proof}
By the Chinese remainder theorem,
$\omega_{420}=\omega_4\cdot\omega_3\cdot\omega_5\cdot\omega_7
= 2\cdot 1\cdot 1\cdot 3 = 6$.
Direct enumeration confirms $\R$ as the six classes
(\texttt{eabc\_hl\_coefficient\_hypotheses.py}).
\end{proof}

\begin{theorem}[Identical HL coefficients on $\R$]
\label{thm:equal-hl}
For all $r\in\R$, $W_{\mathrm{HL}}(r)=W_0$ for the same constant $W_0$.
In particular, the HL heuristic predicts
$P(p\bmod 420=r\mid Q(p)\text{ prime})\to 1/6$.
\end{theorem}

\begin{proof}
Every $r\in\R$ satisfies:
$r\equiv 1\bmod 2$ ($\omega_2=1$);
$r\equiv 2\bmod 3$ ($\omega_3=1$);
$r\equiv 1$ or $3\bmod 4$, three channels each ($\omega_4=2$);
$r\equiv 1\bmod 5$ ($\omega_5=1$);
$r\equiv 2,3$ or $4\bmod 7$, two channels each ($\omega_7=3$).
Each factor in \eqref{eq:channel-weight} is therefore identical
across $\R$.
\end{proof}

\subsection{Local factors and goodness of fit}

Table~\ref{tab:hl-coefficients} lists local congruence structure,
observed occupancies at $N=4767$ ($p\le 10^8$), and identical
HL weights.
Small-prime singular factors are
$w_3=7/8$, $w_5=63/64$, $w_7\approx 0.997$.

\begin{table}[ht]
\centering
\caption{Local HL structure and channel weights ($N=4767$, $p\le 10^8$;
data from \texttt{eabc\_hl\_coefficient\_hypotheses.py}).}
\label{tab:hl-coefficients}
\begin{tabular}{@{}rrrrrrrrr@{}}
\toprule
$r$ & $n_r$ & $\hat p_r$ & $\delta_r$ &
$r\bmod 4$ & $r\bmod 5$ & $r\bmod 7$ &
$W_{\mathrm{HL}}(r)$ & $W_{\mathrm{HL}}/\sum$ \\
\midrule
 11 & 765 & 0.1605 & $-0.00619$ & 3 & 1 & 4 & $3.47\times 10^{-4}$ & $1/6$ \\
101 & 831 & 0.1743 & $+0.00766$ & 1 & 1 & 3 & $3.47\times 10^{-4}$ & $1/6$ \\
191 & 786 & 0.1649 & $-0.00178$ & 3 & 1 & 2 & $3.47\times 10^{-4}$ & $1/6$ \\
221 & 809 & 0.1697 & $+0.00304$ & 1 & 1 & 4 & $3.47\times 10^{-4}$ & $1/6$ \\
311 & 809 & 0.1697 & $+0.00304$ & 3 & 1 & 3 & $3.47\times 10^{-4}$ & $1/6$ \\
401 & 767 & 0.1609 & $-0.00577$ & 1 & 1 & 2 & $3.47\times 10^{-4}$ & $1/6$ \\
\midrule
\multicolumn{8}{l}{Local moduli: $\omega_2{=}1$, $\omega_3{=}1$, $\omega_4{=}2$,
$\omega_5{=}1$, $\omega_7{=}3$.} \\
\bottomrule
\end{tabular}
\end{table}

\begin{table}[ht]
\centering
\caption{HL prediction vs.\ observation ($N=4767$).}
\label{tab:hl-goodness}
\begin{tabular}{@{}lr@{}}
\toprule
Quantity & Value \\
\midrule
$\chi^2$ (HL uniform $1/6$) & $4.344$ \\
$p$-value (5 df) & $0.501$ \\
$\mathrm{KL}(\hat{\mathbf p}\,\|\,\mathbf p_{\mathrm{HL}})$ &
$4.55\times 10^{-4}$ \\
$\max_i |\delta_i|$ & $0.77\%$ \\
$|\langle\delta,\chi_3\rangle|$ & $0.00275$ \\
$|\langle\delta,\chi_3\rangle|/(1/(6\sqrt{N}))$ & $\approx 1.1$ \\
\bottomrule
\end{tabular}
\end{table}

Since all HL weights coincide, $\mathbf p_{\mathrm{HL}}$ and the uniform
prediction $1/6$ are identical; a channel-specific HL correction does
not improve the fit.

\subsection{Hypotheses H1--H6}

We summarize six testable hypotheses
(full analysis in \texttt{eabc\_hl\_hypotheses.tex}):

\begin{description}
  \item[H1 (identical HL coefficients).]
        All six channels carry the same singular-series weight
        (Theorem~\ref{thm:equal-hl}); deviations $\le 0.77\%$.
        \emph{Supported.}
  \item[H2 (mod-$7$ HL densities).]
        Different $p\bmod 7$ classes induce distinct HL factors and
        explain the $\chi_3$ component.
        \emph{Refuted:} $\omega_7=3$ yields equal factors $1/3$ per class;
        $\chi_3|_{\R}=\psi_2$ is algebraic, not density-theoretic.
  \item[H3 (mod-$4$/mod-$5$ HL densities).]
        Alternating $r\bmod 4$ or constant $r\bmod 5$ induce distinct HL
        weights feeding the $\chi_4$ projection.
        \emph{Refuted:} $\omega_4=2$, $\omega_5=1$ give identical factors
        on $\R$.
  \item[H4 (inter-quadruplet correlations).]
        Successive quadruplets are not i.i.d.; transition dynamics modify
        the channel distribution.
        \emph{Open, but unlikely} as the main cause: scaling to $10^9$
        shows no stable drift beyond $O(1/\sqrt{N})$.
  \item[H5 (finite-sample fluctuations).]
        The observed asymmetry at $N=4767$ is pure finite-sample noise
        around HL equidistribution.
        \emph{Supported:} $\chi^2\approx 4.34$, $p\approx 0.50$,
        consistent with Hypothesis~\ref{hyp:H}.
  \item[H6 (unproved analytic equidistribution).]
        Even with identical coefficients, arithmetic correlations may
        distort the conditional distribution $p\bmod 420\mid Q(p)$.
        \emph{Mathematically open}, numerically unremarkable to $10^9$.
\end{description}

%=============================================================================
\section{Discrete Holonomy Interpretation}
\label{sec:holonomy}
%=============================================================================

This section belongs to \textbf{Level~2} of the framework
(Section~\ref{sec:three-levels}): a numerically verified \emph{structural
parallel} between the local $C_3$-phase pattern induced by $\chi_3$ on $\R$
and a discrete edge-phase model on the six-channel ring.
It does not extend the unconditional algebra of Level~1, nor does it assert
any Level~3 physical mechanism.

\subsection{Discrete connection on the six-channel ring}

Fix the channel order $(r_n)_{n=0}^{5}$ from Section~2.
Define vertex phases
\[
  \theta_n := \arg\bigl(\chi_3(r_n)\bigr)
  \in \bigl\{0,\,\tfrac{2\pi}{3},\,-\tfrac{2\pi}{3}\bigr\},
\]
with $\arg$ mapped to the principal value in $(-\pi,\pi]$ after rounding to the
$C_3$ phase set (as implemented in \texttt{eabc\_chi3\_holonomy\_test.py}).
The \emph{discrete connection} along directed edges $n\to n{+}1$ (indices mod~$6$)
is
\[
  \Delta\theta_n := \mathrm{wrap}\bigl(\theta_{n+1}-\theta_n\bigr),
  \qquad
  \mathrm{wrap}(\alpha)\in(-\pi,\pi].
\]
Global holonomy around the closed ring is
\[
  H_\chi := \exp\!\Bigl(i\sum_{n=0}^{5}\Delta\theta_n\Bigr),
  \qquad
  \gamma_\chi := \arg(H_\chi),
  \qquad
  \mathrm{Flux}_\chi := \frac{\gamma_\chi}{2\pi}.
\]
The same construction applies to the $k{=}2$ Fourier mode
$\psi_2(n)=e^{-2\pi i n/3}$ with phases $\phi_n:=\arg(\psi_2(n))$ and edge
increments $\Delta\phi_n:=\mathrm{wrap}(\phi_{n+1}-\phi_n)$.

\subsection{Numerical results}

Direct computation (\texttt{eabc\_chi3\_holonomy\_test.py};
output in \texttt{eabc\_chi3\_holonomy\_report.txt}) yields:

\begin{itemize}
  \item \textbf{Local edge phases ($\chi_3$).}
        $\Delta\theta_n = +2\pi/3$ on all six edges---a constant $C_3$-typical
        step.
        The first half-cycle $(n=0,1,2)$ matches the second $(n=3,4,5)$.
  \item \textbf{Local edge phases ($k{=}2$).}
        $\Delta\phi_n = -2\pi/3$ on all six edges, the connection-level
        structural parallel of Lemma~\ref{lem:chi3-equals-psi2}.
  \item \textbf{Global holonomy.}
        $\gamma_\chi \approx 0$, $\mathrm{Flux}_\chi \approx 0$
        (numerically $\lesssim 10^{-15}$);
        likewise $\gamma_{\psi_2}\approx 0$, $\mathrm{Flux}_{\psi_2}\approx 0$.
        The wrapped edge increments sum to an integer multiple of $2\pi$, so
        the closed six-ring carries trivial total phase holonomy---as expected
        for a defect-free cycle with no enclosed flux in this discrete model.
  \item \textbf{Null test ($\chi_5$).}
        Since $\chi_5(r)=+1$ for all $r\in\R$, $\theta_n\equiv 0$ and all
        $\Delta\theta_n=0$; holonomy is identically trivial.
\end{itemize}

These findings confirm a \emph{coordinate-theoretic analogy}: the mod-$7$ cubic
character induces a uniform local phase-transport pattern on the channel ring,
while global holonomy cancels on the closed cycle.
They mirror the exact identity $\chi_3|_{\R}=\psi_2$ (Level~1) in a
graph-theoretic language of edge phases.

\paragraph{Epistemological disclaimer.}
This establishes a coordinate-theoretic analogy, not a claim about quantized
transport or persistent topological bias in prime quadruplets.
We do \emph{not} assert quantum Hall plateaus, topological protection in the
sense of the von~Klitzing effect, an Aharonov--Bohm solenoid experiment, or
non-abelian gauge fields realized in nature.
The discrete edge-phase model is a structural parallel useful for organizing
$\chi_3$-data on $\R$; no physical mechanism coupling holonomy to quadruplet
occupancy has been identified or tested.

%=============================================================================
\section{Conclusion}
%=============================================================================

The modulo-$420$ channel asymmetry of prime quadruplets admits an exact
character-theoretic diagonalization.
The dominant Fourier component coincides with the restriction of the cubic
Dirichlet character modulo $7$, while the remaining structure is accounted
for by quadratic character components.
However, scaling experiments up to $10^9$ show no evidence for a persistent
Chebyshev-type drift.
The observed amplitudes remain compatible with multinomial fluctuations
around the uniform distribution.
Consequently, the present data support a character-space decomposition of
the fluctuations rather than a genuine asymptotic bias.

A detailed Hardy--Littlewood analysis of the singular series shows
identical coefficients for all six modulo-$420$ channels.
Asymptotically, the Hardy--Littlewood heuristic therefore predicts strict
equidistribution of prime quadruplets across the channels.
The empirically observed asymmetry vector is statistically consistent with
pure finite-sample fluctuations ($\chi^2 \approx 4.34$, $p \approx 0.50$).
The previously established mathematical dominance of the cubic character
$\chi_3$ is consequently not density-theoretic but purely algebraically
enforced by the reduction of the underlying residue-class space.

The remaining open problem is an analytic equidistribution theorem for
prime quadruplets on the six admissible modulo-$420$ channels.

Section~\ref{sec:holonomy} adds a Level~2 discrete holonomy interpretation:
the local $C_3$ edge-phase pattern induced by $\chi_3$ is numerically
consistent with a trivial global holonomy on the closed six-ring, furnishing
a coordinate-theoretic analogy only---not evidence for quantized transport
or persistent topological bias.

\paragraph{Precision.}
The diagonalization refers to the \emph{observed finite sample} on $\R$:
$\chi_3$ supplies a fixed character coordinate system, not a mechanism whose
coefficients grow with $N$.
At $N_{\max}=10^9$ ($K=28387$), $|\langle\delta,\chi_3\rangle|$ scales as
$1/\sqrt{K}$ with $z\approx 1$, consistent with $f(N)=1/(6\sqrt{N})$ under $H$
and inconsistent with a persistent structural bias.

%=============================================================================
\section{Epistemological Scope}
%=============================================================================

The analysis follows a three-stage epistemic chain:
\emph{phenomenon} (observed channel asymmetry)
$\to$ \emph{diagonalization} (exact character decomposition, Level~1)
$\to$ \emph{discrete holonomy analogy} (numerical edge-phase model,
Section~\ref{sec:holonomy}, Level~2)
$\to$ \emph{HL corroboration} (singular-series weights, Level~3).
Algebra and asymptotics are deliberately decoupled:
Level~1 statements hold unconditionally on the fixed finite set $\R$,
while Level~2 and Level~3 address the \emph{magnitudes} of coefficients
under probabilistic or arithmetic hypotheses.

\begin{itemize}
  \item \textbf{Confirmed (algebraic diagonalization).}
        The EABC channel labels organize a finite residue set on which
        $\chi_3|_{\R}=\psi_2$ holds exactly.
        Decomposing $\delta$ in Dirichlet characters is a proved descriptive
        tool for the observed fluctuations.
        The $\chi_3$ dominance is enforced by the residue-class geometry,
        not by unequal HL densities (H2, H3 refuted).

  \item \textbf{Confirmed (discrete holonomy analogy, Level~2).}
        On the six-channel ring, $\chi_3$ induces constant local edge phases
        $\Delta\theta_n=+2\pi/3$ with trivial global holonomy
        ($\gamma_\chi\approx 0$, $\mathrm{Flux}_\chi\approx 0$);
        the $\chi_5$ null test is trivial.
        This is a structural parallel between character phases and a discrete
        edge-phase model---not a claim about topological protection,
        quantum Hall plateaus, or physical gauge flux
        (Section~\ref{sec:holonomy}).

  \item \textbf{Not confirmed (new asymptotic bias).}
        Neither $B_{101,11}/\sqrt{N}$ nor $|\langle\delta_N,\chi_3\rangle|$
        exhibits growth beyond $O_p(1/\sqrt{N})$.
        A channel-specific Chebyshev bias is not supported by data to $10^9$.

  \item \textbf{Corroborated (HL singular-series structure).}
        Theorem~\ref{thm:equal-hl} shows identical HL coefficients on all
        six channels; the observed occupancy vector is consistent with
        uniform $1/6$ ($\chi^2\approx 4.34$, $p\approx 0.50$,
        $\mathrm{KL}\approx 4.55\times 10^{-4}$).
        This arithmetic input justifies Hypothesis~$H$ as the natural
        asymptotic model, without claiming a completed proof.

  \item \textbf{Open (analytic equidistribution, H6).}
        Proving conditional equidistribution of $p\bmod 420$ given
        $Q(p)$ prime, and ruling out correlations beyond $O(1/\sqrt{N})$,
        remains a separate analytic problem.
\end{itemize}

Diagonalization answers \emph{in which directions} the observed bias
decomposes; HL analysis answers \emph{whether unequal arithmetic densities}
could account for those magnitudes (they cannot on $\R$); an analytic
equidistribution theorem would close the remaining proof gap.

%=============================================================================
\begin{thebibliography}{9}

\bibitem{chebyshev1853}
P.~L.~Chebyshev,
\emph{Sur la concurrence des s\'eries trigonom\'etriques},
M\'em.\ Acad.\ Imp.\ St.-P\'etersbourg \textbf{3} (1853).

\bibitem{dirichlet1837}
P.~G.~L.~Dirichlet,
\emph{Beweis des Satzes, dass jede unbegrenzte arithmetische Progression,
deren erstes Glied und Differenz ganze Zahlen ohne gemeinschaftlichen
Factor sind, unendlich viele Primzahlen enth\"alt},
Abh.\ K\"onigl.\ Preu\ss.\ Akad.\ Wiss.\ (1837).

\bibitem{rubinstein2001}
M.~Rubinstein and P.~Sarnak,
\emph{Chebyshev's bias},
Experiment.\ Math.\ \textbf{3} (1994), 173--197.

\bibitem{hardylittlewood1923}
G.~H.~Hardy and J.~E.~Littlewood,
\emph{Some problems of `Partitio numerorum' III:
On the expression of a number as a sum of primes},
Acta Math.\ \textbf{44} (1923), 1--70.

\bibitem{batemanhorn1962}
P.~T.~Bateman and R.~A.~Horn,
\emph{Primes in arithmetic progressions and linear forms of primes},
in \emph{Proc.\ Fourth Canadian Math.\ Congress},
Univ.\ Toronto Press, 1962, 193--201.

\bibitem{davenport2000}
H.~Davenport,
\emph{Multiplicative Number Theory},
3rd ed., Springer, 2000.

\end{thebibliography}

%=============================================================================
\appendix
\section{Computational Reproducibility}
%=============================================================================

The numerical claims are reproducible from the following scripts in the
same directory:

\begin{itemize}
  \item \texttt{eabc\_spectral\_decomposition.py} --- character decomposition
        of $\delta$ at $K=4767$; verifies $\chi_3|_{\R}=\psi_2$ and
        Tables~\ref{tab:projections}--\ref{tab:char-values}.
  \item \texttt{eabc\_scaling\_test.py} --- pair bias $B_{101,11}$ and
        $B/\sqrt{K}$ across $10^8$, $3\cdot 10^8$, $10^9$
        (Table~\ref{tab:B-scaling}).
  \item \texttt{eabc\_analytic\_proof\_check.py} --- $\chi_3$-projection
        scaling and multinomial null $p$-values (Table~\ref{tab:chi3-scaling}).
  \item \texttt{eabc\_chi3\_mod7\_projection.py} --- mod-$7$ cubic projection
        diagnostics.
  \item \texttt{eabc\_dirichlet\_mod420\_projection.py} --- full mod-$420$
        Dirichlet projection machinery.
  \item \texttt{eabc\_cumulative\_bias.py} --- cumulative bias trajectories.
  \item \texttt{eabc\_hl\_coefficient\_hypotheses.py} --- Hardy--Littlewood
        singular-series computation, local factors mod $2,3,4,5,7$,
        goodness-of-fit statistics (Tables~\ref{tab:hl-coefficients}--\ref{tab:hl-goodness};
        output in \texttt{eabc\_hl\_coefficient\_hypotheses.json} and
        \texttt{eabc\_hl\_coefficient\_hypotheses\_report.txt}).
  \item \texttt{eabc\_hl\_hypotheses.tex} --- standalone HL hypothesis
        analysis (H1--H6) in German.
  \item \texttt{eabc\_chi3\_holonomy\_test.py} --- discrete edge-phase
        holonomy on the six-channel ring for $\chi_3$, $k{=}2$ Fourier mode,
        and $\chi_5$ null test (Section~\ref{sec:holonomy};
        output in \texttt{eabc\_chi3\_holonomy\_report.txt}).
\end{itemize}

\end{document}
